\label{sec:benchmarks}

The benchmark suite is intentionally heterogeneous. The central claim of the
paper is not that one solver policy dominates every nonlinear PDE family, but
that one maintained differentiable sparse-solver stack can cover scalar
elliptic problems, phase-field-type energies, finite-strain mechanics,
elastoplastic slope stability, and reduced topology optimization without
switching conceptual frameworks.

\begin{table}[H]
  \centering
  \begin{tabular}{@{}l c >{\RaggedRight\arraybackslash}p{0.14\textwidth} >{\RaggedRight\arraybackslash}p{0.23\textwidth} >{\RaggedRight\arraybackslash}p{0.26\textwidth}@{}}
\toprule
Family & Grid / mesh & Stop rule & Compared paths & Main difficulty \\
\midrule
$p$-Laplace & $L_{9}$ & Newton + line search & FEniCS, pure JAX, JAX+PETSc & nonlinear elliptic solve with exact sparse Hessians \\
Ginzburg--Landau & $L_{9}$ & Newton + line search & FEniCS, JAX+PETSc & indefinite local curvature from the double well \\
Hyperelasticity & $L_{4}$, 24 steps & trust-region path & FEniCS, pure JAX, JAX+PETSc & nonconvex large-deformation mechanics \\
Plasticity2D & $P_{4}(L_{5})$--$P_{4}(L_{7})$ & endpoint or fixed work & JAX+PETSc only & same-mesh PMG and nonlinear tail behavior \\
Plasticity3D & $P_{1}$/$P_{2}$/$P_{4}$ & $\|g\| < \num{0.01}$ & constitutive-AD and source PMG variants & heterogeneous 3D Mohr--Coulomb with constitutive AD \\
Topology & $768\times384$ & stall-stop continuation & pure JAX, JAX+PETSc & distributed design-mechanics coupling \\
\bottomrule
\end{tabular}

  \caption{Benchmark families, stopping rules, and principal comparison paths.
  The later numerical section revisits the same families from the timing and
  scaling viewpoint.}
  \label{tab:benchmark-spec}
\end{table}

\begin{table}[H]
  \centering
  \begin{tabularx}{\textwidth}{@{}>{\hsize=1.080\hsize\linewidth=\hsize\RaggedRight\arraybackslash}X@{\hspace{0.8em}}>{\hsize=0.580\hsize\linewidth=\hsize\Centering\arraybackslash}X@{\hspace{0.7em}}>{\hsize=0.620\hsize\linewidth=\hsize\Centering\arraybackslash}X@{\hspace{1.0em}}>{\hsize=1.720\hsize\linewidth=\hsize\RaggedRight\arraybackslash}X@{}}
\toprule
Family & FEniCS & pure JAX & Notes \\
\midrule
$p$-Laplace & yes & yes & All three stacks exist on maintained showcase cases. \\
Ginzburg--Landau & yes & no & FEniCS and JAX+PETSc form the maintained comparison. \\
Hyperelasticity & yes & yes & pure JAX is a serial formulation reference only. \\
Plasticity2D & no & no & The maintained story is JAX+PETSc only. \\
Plasticity3D & no & no & Source assembly exists as supporting comparison, not as a maintained reference path. \\
Topology & no & yes & Parallel fine-grid path is JAX+PETSc; pure JAX remains the serial design reference. \\
\bottomrule
\end{tabularx}

  \caption{Availability of \fenics{} and \purejax{} references. Reference
  paths are used where the repository already provides documented implementations;
  absent cells are left absent rather than filled with speculative baselines.}
  \label{tab:reference-availability}
\end{table}

\subsection{\texorpdfstring{$p$-Laplace}{p-Laplace}}
\label{sec:bench-plaplace}

The literature model is the standard stationary $p$-Laplace Dirichlet problem
and its associated energy formulation \citep[Ch.~2]{lindqvist2019plaplace}. The
unknown is a scalar field $u:\Omega \to \mathbb{R}$ and the continuous energy is
\begin{equation}
  \mathcal{E}(u)
  = \int_{\Omega} \frac{1}{p} \lvert \nabla u \rvert^p \, dx
    - \int_{\Omega} f u \, dx,
  \label{eq:plaplace-energy}
\end{equation}
whose Euler--Lagrange equation is
\begin{equation}
  -\nabla \cdot \left(\lvert \nabla u \rvert^{p-2} \nabla u \right) = f
  \quad \text{in }\Omega,
  \qquad
  u = 0 \quad \text{on }\partial\Omega.
  \label{eq:plaplace-strong}
\end{equation}
The weak problem is: find $u \in W_0^{1,p}(\Omega)$ such that
\begin{equation}
  \int_{\Omega}
    \lvert \nabla u \rvert^{p-2} \nabla u \cdot \nabla v \, dx
  =
  \int_{\Omega} f v \, dx
  \qquad \forall v \in W_0^{1,p}(\Omega),
  \label{eq:plaplace-weak}
\end{equation}
The repository specialization used in this paper takes
$\Omega = (0,1)^2$, exponent $p=3$, forcing $f=-10$, and a conforming $P_1$
triangular hierarchy. The discrete problem is
\begin{equation}
  u_h \in V_h \subset W_0^{1,p}(\Omega):
  \qquad
  u_h = \arg\min_{v_h \in V_h} \mathcal{E}_h(v_h),
  \label{eq:plaplace-discrete}
\end{equation}
where $V_h$ is the chosen conforming finite-element subspace. This is the
cleanest benchmark in the paper for checking exact gradients, exact Hessians,
and sparse-recovery alternatives against each other without constitutive or
geometric complications.

\begin{figure}[H]
  \centering
  \includegraphics[width=\PaperFigureNarrow]{plaplace_state.pdf}
  \caption{Representative $p$-Laplace state on the maintained hierarchy.}
  \label{fig:plaplace-state}
\end{figure}

\begin{figure}[H]
  \centering
  \includegraphics[width=\PaperFigureNarrow]{plaplace_energy_levels.pdf}
  \caption{Final discrete energies from \cref{eq:plaplace-discrete} across the
  maintained mesh hierarchy. The element-AD and local-SFD \jaxpetsc{} paths are
  essentially energy-consistent with the \fenics{} reference.}
  \label{fig:plaplace-energy-levels}
\end{figure}

\begin{table}[H]
  \centering
  \begin{tabular*}{\textwidth}{@{}l@{\extracolsep{\fill}} c c c c@{}}
\toprule
Path & Energy & Newton iters & Krylov iters & Wall time [s] \\
\midrule
FEniCS custom Newton & \num{-7.942968719} & 5 & 15 & \num{0.0782} \\
JAX+PETSc element AD & \num{-7.942968719} & 6 & 17 & \num{0.0798} \\
JAX+PETSc colored SFD & \num{-7.942968711} & 6 & 6 & \num{0.194} \\
\bottomrule
\end{tabular*}

  \caption{Representative final discrete energies from
  \cref{eq:plaplace-discrete} on the showcase case.}
  \label{tab:plaplace-benchmark}
\end{table}

The benchmark-local picture is therefore one of formulation agreement:
\cref{fig:plaplace-energy-levels,tab:plaplace-benchmark} show that the
maintained routes converge to the same discrete energy from
\cref{eq:plaplace-discrete}. Scaling and timing for this family are reported in
\Cref{sec:num-plaplace}.

\subsection{Ginzburg--Landau}
\label{sec:bench-gl}

We cite \citet{ginzburg1950theory} for the historical origin of
Ginzburg--Landau free-energy modeling only. The maintained benchmark used here is a
repository specialization: a scalar real-valued double-well problem on
$\Omega = [-1,1]^2$ with homogeneous Dirichlet boundary conditions and energy
\begin{equation}
  \mathcal{E}(u)
  =
  \int_{\Omega}
    \frac{\varepsilon}{2} \lvert \nabla u \rvert^2
    + \frac{1}{4}(u^2 - 1)^2 \, dx,
  \qquad
  \varepsilon = 10^{-2}.
  \label{eq:gl-energy}
\end{equation}
Its strong form is
\begin{equation}
  -\varepsilon \Delta u + u(u^2 - 1) = 0
  \quad \text{in }\Omega,
  \qquad
  u = 0 \quad \text{on }\partial\Omega,
  \label{eq:gl-strong}
\end{equation}
and the weak problem is: find $u \in H_0^1(\Omega)$ such that
\begin{equation}
  \int_{\Omega}
    \varepsilon \nabla u \cdot \nabla v
    + (u^2 - 1) u v \, dx
  = 0
  \qquad \forall v \in H_0^1(\Omega).
  \label{eq:gl-weak}
\end{equation}
The maintained discrete problem is a conforming $P_1$ stationary problem on the
same triangular hierarchy:
\begin{equation}
  u_h \in V_h:
  \qquad
  D\mathcal{E}_h(u_h)[v_h] = 0
  \qquad \forall v_h \in V_h.
  \label{eq:gl-discrete}
\end{equation}
Relative to $p$-Laplace, the additional difficulty here is not geometry but
indefinite local curvature from the double-well potential in
\cref{eq:gl-energy}.

\begin{figure}[H]
  \centering
  \includegraphics[width=\PaperFigureNarrow]{ginzburg_landau_state.pdf}
  \caption{Representative Ginzburg--Landau state on the maintained hierarchy.}
  \label{fig:gl-state}
\end{figure}

\begin{figure}[H]
  \centering
  \includegraphics[width=\PaperFigureNarrow]{ginzburg_landau_energy_levels.pdf}
  \caption{Final discrete energies associated with the stationary problem
  \cref{eq:gl-discrete} across mesh
  levels. The maintained element-AD and local-SFD routes remain nearly
  coincident with the \fenics{} reference.}
  \label{fig:gl-energy-levels}
\end{figure}

\begin{table}[H]
  \centering
  \begin{tabular*}{\textwidth}{@{}l@{\extracolsep{\fill}} c c c c@{}}
\toprule
Path & Energy & Newton iters & Krylov iters & Wall time [s] \\
\midrule
FEniCS custom Newton & \num{0.3462323651} & 8 & 33 & \num{0.113} \\
JAX+PETSc element AD & \num{0.3462314438} & 8 & 33 & \num{0.0949} \\
JAX+PETSc colored SFD & \num{0.3462314438} & 8 & 33 & \num{0.247} \\
\bottomrule
\end{tabular*}

  \caption{Representative final discrete energies associated with
  \cref{eq:gl-discrete} on the showcase case.}
  \label{tab:gl-benchmark}
\end{table}

The local picture is again one of formulation consistency, now for a more
solver-sensitive nonconvex energy. Strong scaling and solver-time comparisons
for this family are reported in \Cref{sec:num-gl}.

\subsection{Hyperelasticity}
\label{sec:bench-he}

The literature model is a standard finite-strain hyperelastic equilibrium
problem written in terms of deformation map $y(X) = X + u(X)$, deformation
gradient $F = \nabla y$, Jacobian $J = \det F$, and first Piola stress
$P = \partial W / \partial F$ \citep[Chs.~4--6]{bonet2008nonlinear}. The
repository specialization used in this paper is a cantilever benchmark with the
compressible Neo-Hookean-type density
\begin{equation}
  \Pi(y;\eta)
  =
  \int_{\Omega}
    C_1 \left( \operatorname{tr}(F^T F) - 3 - 2 \ln J \right)
    + D_1 (J - 1)^2 \, dx,
  \label{eq:hyperelasticity-energy}
\end{equation}
on $\Omega = [0,0.4] \times [-0.005,0.005]^2$ with
$C_1 = \num{3.8461538461538464e7}$ and
$D_1 = \num{8.333333333333333e7}$. The left face is clamped, while the right
face carries the prescribed rotating Dirichlet load path used throughout the
repository, parametrized by load step $\eta$. The strong form is
\begin{equation}
  -\nabla \cdot P(F) = 0 \quad \text{in }\Omega,
  \qquad
  y = \bar y(\eta) \quad \text{on }\Gamma_D,
  \qquad
  P(F) n = 0 \quad \text{on }\Gamma_N,
  \label{eq:hyperelasticity-strong}
\end{equation}
where $P = \partial W / \partial F$ is the first Piola stress. The weak form is:
find $u \in V_{\bar y(\eta)}$ such that
\begin{equation}
  \int_{\Omega} P(F(u)) : \nabla v \, dx = 0
  \qquad \forall v \in V_0.
  \label{eq:hyperelasticity-weak}
\end{equation}
The discrete problem solved in the paper is the sequence of stationary points
\begin{equation}
  u_h^{(k)} \in V_h:
  \qquad
  D\Pi_h(u_h^{(k)};\eta_k)[v_h] = 0
  \qquad \forall v_h \in V_{h,0},
  \qquad
  k = 1,\dots,24,
  \label{eq:hyperelasticity-discrete}
\end{equation}
with first-order tetrahedral vector elements and the benchmark-specific boundary
data described above. This is the first family in which globalization is more
scientifically important than derivative parity alone.

\begin{figure}[H]
  \centering
  \includegraphics[width=\PaperFigureMedium]{hyperelasticity_state.pdf}
  \caption{Representative finite-strain hyperelastic state, colored by the
  energy density in \cref{eq:hyperelasticity-energy}.}
  \label{fig:hyperelasticity-state}
\end{figure}

\begin{figure}[H]
  \centering
  \includegraphics[width=\PaperFigureNarrow]{hyperelasticity_energy_levels.pdf}
  \caption{Final discrete energies along the stationary sequence
  \cref{eq:hyperelasticity-discrete} across the maintained mesh hierarchy. The
  publication story is less about
  exact equality than about tracking the same 24-step trajectory consistently.}
  \label{fig:hyperelasticity-energy-levels}
\end{figure}

\begin{table}[H]
  \centering
  \begin{tabular*}{\textwidth}{@{}l@{\extracolsep{\fill}} c c c c@{}}
\toprule
Path & Energy & Steps & Krylov iters & Wall time [s] \\
\midrule
FEniCS custom Newton & \num{197.775074} & 24 & 8722 & \num{6.51} \\
JAX+PETSc element AD & \num{197.754582} & 24 & 10595 & \num{9.36} \\
serial JAX & \num{197.749509} & 24 & 2284 & \num{41.5} \\
\bottomrule
\end{tabular*}

  \caption{Representative 24-step terminal energies along
  \cref{eq:hyperelasticity-discrete}.}
  \label{tab:hyperelasticity-benchmark}
\end{table}

The benchmark-local takeaway is therefore a globalization result:
\cref{fig:hyperelasticity-state,fig:hyperelasticity-energy-levels} show the
terminal deformed state and terminal energies, while the scaling evidence is
deferred to \Cref{sec:num-hyperelasticity}. That later section also adds one
matched serial companion baseline against JAX-FEM on the same mesh, material
law, and displacement schedule; it is included as a fairness-first
implementation comparison rather than as a broad framework ranking exercise.

\subsection{Plasticity2D}
\label{sec:bench-plasticity2d}

The literature model is a quasi-static plane-strain Mohr--Coulomb slope
stability problem posed as equilibrium together with a constitutive update and
history variables
\citep{davis1968plasticity,simo1985consistent,simo1998compinel,tschuchnigg2015nonassociated,sysala2017returnmapping,sysala2021optimization}.
The unknown is the displacement field
$u = (u_x,u_y):\Omega \to \mathbb{R}^2$ on the maintained homogeneous slope
geometry. At increment $n+1$, the continuum equilibrium equations are
\begin{equation}
  -\nabla \cdot \sigma^{n+1} = b^{n+1}
  \quad \text{in }\Omega,
  \qquad
  u^{n+1} = \bar u^{n+1} \quad \text{on }\Gamma_D,
  \qquad
  \sigma^{n+1}n = t^{n+1} \quad \text{on }\Gamma_N,
  \label{eq:plasticity2d-strong}
\end{equation}
where $\sigma^{n+1}$ and the internal variables are obtained from the
elastoplastic update driven by $\varepsilon(u^{n+1})$ and the previous state.
The corresponding weak form is
\begin{equation}
  \int_{\Omega} \sigma^{n+1} : \varepsilon(v) \, dx
  =
  \int_{\Omega} b^{n+1} \cdot v \, dx
  + \int_{\Gamma_N} t^{n+1} \cdot v \, ds
  \qquad \forall v \in V_0.
  \label{eq:plasticity2d-weak}
\end{equation}
The maintained runs reported here are endpoint solves, not path-consistent
incremental histories. Accordingly, the repository benchmark uses the following
scalarized discrete surrogate functional for differentiation and comparison:
\begin{equation}
  \Pi_h^{2\mathrm{D}}(u_h;\lambda_{\mathrm{sr}})
  =
  \sum_{e=1}^{n_{\mathrm{el}}}
  \sum_{q=1}^{n_q}
    w_{eq}\,
    \Phi_{\mathrm{MC},2\mathrm{D}}
    \!\left(
      B_{eq}u_e,\varepsilon^{p,\mathrm{old}}_{eq},
      E,\nu,c_{\lambda,eq},\phi_{\lambda,eq},\psi_{\lambda,eq}
    \right)
  - f_{\mathrm{ext}}^T u_h,
  \label{eq:plasticity2d-discrete}
\end{equation}
where $\lambda_{\mathrm{sr}}$ denotes the strength-reduction factor, distinct
from the Lam{\'e} parameter used later in three dimensions.

The maintained code keeps a Davis-B-style reduction path to map the raw
cohesion and angles to reduced values under strength reduction; with
$c_0,\phi,\psi$ denoting the unreduced cohesion, friction angle, and dilation
angle, respectively, the benchmark uses
\begin{align}
  c_1 &= c_0 / \lambda_{\mathrm{sr}},
  &
  \phi_1 &= \arctan \!\left(\tan \phi / \lambda_{\mathrm{sr}}\right),
  &
  \psi_1 &= \arctan \!\left(\tan \psi / \lambda_{\mathrm{sr}}\right),
  \nonumber\\
  \beta &= \frac{\cos \phi_1 \cos \psi_1}{1 - \sin \phi_1 \sin \psi_1},
  &
  c_{\lambda} &= \beta c_1,
  &
  \phi_{\lambda} &= \arctan\!\left(\beta \tan \phi_1\right),
  \label{eq:plasticity2d-davis}
\end{align}
with $\psi_{\lambda}$ defined analogously when the non-associated flow rule is
activated \citep{davis1968plasticity,tschuchnigg2015nonassociated,sysala2021optimization}. The
showcase endpoint runs in this paper use $\lambda_{\mathrm{sr}} = 1$ and
$\psi = 0^\circ$, so the reported states should be read as repository endpoint
surrogates rather than path-consistent history updates.

\begin{figure}[H]
  \centering
  \includegraphics[width=\PaperFigureMedium]{plasticity2d_state_pair.pdf}
  \caption{Representative Plasticity2D state: displacement magnitude and
  deviatoric strain derived from the scalar surrogate
  \cref{eq:plasticity2d-discrete}.}
  \label{fig:plasticity2d-state}
\end{figure}

\begin{figure}[H]
  \centering
  \includegraphics[width=\PaperFigureNarrow]{plasticity2d_resolution_energy.pdf}
  \caption{Resolution trend for the maintained Plasticity2D family. The plotted
  terminal energies are evaluations of \cref{eq:plasticity2d-discrete} on the
  curated $P4(L5)$ showcase and on the larger fixed-work $P4(L6)$ and
  $P4(L7)$ campaigns.}
  \label{fig:plasticity2d-energy-levels}
\end{figure}

\begin{table}[H]
  \centering
  \begin{tabular}{@{}l c c c l@{}}
\toprule
Case & Free DOFs & Energy & Wall time [s] & Note \\
\midrule
$P_{4}(L_{5})$ & \num{332960} & \num{-212.538} & \num{147} & curated showcase \\
$P_{4}(L_{6})$ & \num{1331520} & \num{-212.539} & \num{147} & fixed-work at 8 ranks \\
$P_{4}(L_{7})$ & \num{5325440} & \num{-212.539} & \num{253} & fixed-work at 16 ranks \\
\bottomrule
\end{tabular}

  \caption{Representative Plasticity2D endpoint values from
  \cref{eq:plasticity2d-discrete}.}
  \label{tab:plasticity2d-benchmark}
\end{table}

This family is where multigrid hierarchy design first becomes dominant in the
paper. Solver-engineering evidence for Plasticity2D is reported later in
\Cref{sec:num-plasticity2d,sec:solver-evidence}.

\subsection{Plasticity3D}
\label{sec:bench-plasticity3d}

The literature model is again quasi-static Mohr--Coulomb equilibrium with
constitutive/history updates, now on a three-dimensional heterogeneous slope
geometry
\citep{davis1968plasticity,simo1985consistent,simo1998compinel,tschuchnigg2015nonassociated,sysala2017returnmapping,sysala2025advancedcontinuation}.
The unknown is a displacement field
$u = (u_x,u_y,u_z): \Omega \to \mathbb{R}^3$, gravity acts in the source
$y$-direction. At increment $n+1$, the continuum equilibrium equations are
\begin{equation}
  -\nabla \cdot \sigma^{n+1} = \rho g_y^{n+1} e_y
  \quad \text{in }\Omega,
  \qquad
  u_i^{n+1} = \bar u_i^{n+1} \quad \text{on }\Gamma_{D,i},
  \qquad
  \sigma^{n+1}n = 0 \quad \text{on }\Gamma_N,
  \label{eq:plasticity3d-strong}
\end{equation}
where $\sigma^{n+1}$ and the internal variables are obtained from the
elastoplastic update driven by $\varepsilon(u^{n+1})$ and the previous state.
The corresponding weak form is
\begin{equation}
  \int_{\Omega} \sigma^{n+1} : \varepsilon(v) \, dx
  =
  \int_{\Omega} \rho g_y^{n+1} e_y \cdot v \, dx
  \qquad \forall v \in V_0.
  \label{eq:plasticity3d-weak}
\end{equation}
The repository specialization studied in this paper uses the imported SSR mesh
hierarchy with same-mesh tetrahedral $P1/P2/P4$ spaces, source component-wise
boundary labels, and an additional glued-bottom condition that fixes every node
with $|y| \le 10^{-12}$ in all three displacement components. The benchmark-local
figures below use the corrected glued-bottom $\code{P4(L1\_2)}$ showcase at
$\lambda_{\mathrm{sr}}=1.55$ and the finished nine-case maintained-local study
at the same load factor. The promoted $\lambda_{\mathrm{sr}}=1.0$ strong-scaling
campaign is kept later in the paper as separate historical timing evidence.
For those endpoint and continuation studies, the maintained code evaluates the
following algorithmic scalar surrogate:
\begin{equation}
  \Pi_h^{3\mathrm{D}}(u_h;\lambda_{\mathrm{sr}})
  =
  \sum_{e=1}^{n_{\mathrm{el}}}
  \sum_{q=1}^{n_q}
    w_{eq}\,
    \Phi_{\mathrm{MC},3\mathrm{D}}
    \!\left(
      \varepsilon_{eq}(u_e),
      \bar c_{\lambda,eq},
      \sin \phi_{\lambda,eq},
      \mu_{eq},
      K_{eq},
      \lambda_{\mathrm{L},eq}
    \right)
  - f_{\mathrm{ext}}^T u_h,
  \label{eq:plasticity3d-discrete}
\end{equation}
where $\lambda_{\mathrm{sr}}$ is again the strength-reduction factor and
$\lambda_{\mathrm{L}}$ is the Lam{\'e} parameter of the elastic predictor.

The code uses the engineering-strain ordering
\begin{equation}
  \varepsilon
  =
  \begin{bmatrix}
    \varepsilon_{xx} &
    \varepsilon_{yy} &
    \varepsilon_{zz} &
    \gamma_{xy} &
    \gamma_{yz} &
    \gamma_{xz}
  \end{bmatrix}^{T},
  \label{eq:plasticity3d-engineering-strain}
\end{equation}
which is first mapped to tensor shear before principal values are formed. The
benchmark keeps the Davis-B reduction explicitly at quadrature points:
\begin{align}
  c_{1,q} &= c_{0,q}/\lambda_{\mathrm{sr}},
  &
  \phi_{1,q} &= \arctan\!\left(\tan\phi_q/\lambda_{\mathrm{sr}}\right),
  &
  \psi_{1,q} &= \arctan\!\left(\tan\psi_q/\lambda_{\mathrm{sr}}\right),
  \nonumber\\
  \beta_q &= \frac{\cos\phi_{1,q}\cos\psi_{1,q}}
                 {1-\sin\phi_{1,q}\sin\psi_{1,q}},
  &
  c_{\lambda,q} &= \beta_q c_{1,q},
  &
  \phi_{\lambda,q} &= \arctan\!\left(\beta_q \tan\phi_{1,q}\right),
  \label{eq:plasticity3d-davis}
\end{align}
with
\begin{equation}
  \bar c_{\lambda,q} = 2 c_{\lambda,q}\cos\phi_{\lambda,q},
  \qquad
  \sin\phi_{\lambda,q} = \sin(\phi_{\lambda,q}),
  \label{eq:plasticity3d-davis-bar}
\end{equation}
which are the scalar constitutive parameters actually stored in the maintained
3D assembly path
\citep{davis1968plasticity,tschuchnigg2015nonassociated,sysala2021optimization}.
Here $\phi$ denotes friction angle and $\psi$ denotes dilation angle; the paper
uses those symbols only in that geomechanics sense.

The local constitutive surrogate follows the piecewise branch structure used in
the repository. Let $\varepsilon_1 \ge \varepsilon_2 \ge \varepsilon_3$ be the
principal strains and let
\begin{equation}
  I_1 = \varepsilon_{xx} + \varepsilon_{yy} + \varepsilon_{zz},
  \qquad
  q(\varepsilon) =
  \varepsilon_{xx}^2 + \varepsilon_{yy}^2 + \varepsilon_{zz}^2
  + \frac{1}{2}\left(
    \gamma_{xy}^2 + \gamma_{yz}^2 + \gamma_{xz}^2
  \right).
  \label{eq:plasticity3d-invariants}
\end{equation}
The elastic branch is
\begin{equation}
  \Phi_{\mathrm{el}}
  = \frac{1}{2}\lambda_{\mathrm{L}} I_1^2 + \mu q(\varepsilon),
  \label{eq:plasticity3d-elastic-branch}
\end{equation}
and the source-style return variables are
\begin{align}
  f_{\mathrm{tr}}
  &=
  2\mu\!\left((1+\sin\phi_{\lambda})\varepsilon_1 -
  (1-\sin\phi_{\lambda})\varepsilon_3\right)
  + 2\lambda_{\mathrm{L}}\sin\phi_{\lambda} I_1
  - \bar c_{\lambda},
  \nonumber\\
  \gamma_{\mathrm{sl}}
  &= \frac{\varepsilon_1-\varepsilon_2}{1+\sin\phi_{\lambda}},
  \qquad
  \gamma_{\mathrm{sr}}
  = \frac{\varepsilon_2-\varepsilon_3}{1-\sin\phi_{\lambda}},
  \label{eq:plasticity3d-return-tests}
\end{align}
with the remaining left, right, and apex thresholds defined as in the code
through the same principal-strain ordering. The scalar potential is then the
piecewise function
\begin{equation}
  \Phi_{\mathrm{MC},3\mathrm{D}}
  =
  \begin{cases}
    \Phi_{\mathrm{el}},
      & f_{\mathrm{tr}} \le 0, \\
    \Phi_{\mathrm{s}},
      & \Lambda_{\mathrm{s}} \le \min(\gamma_{\mathrm{sl}},\gamma_{\mathrm{sr}}), \\
    \Phi_{\mathrm{l}},
      & \gamma_{\mathrm{sl}} < \gamma_{\mathrm{sr}}
        \text{ and }
        \gamma_{\mathrm{sl}} \le \Lambda_{\mathrm{l}} \le \gamma_{\mathrm{la}}, \\
    \Phi_{\mathrm{r}},
      & \gamma_{\mathrm{sr}} < \gamma_{\mathrm{sl}}
        \text{ and }
        \gamma_{\mathrm{sr}} \le \Lambda_{\mathrm{r}} \le \gamma_{\mathrm{ra}}, \\
    \Phi_{\mathrm{a}},
      & \text{otherwise},
  \end{cases}
  \label{eq:plasticity3d-potential}
\end{equation}
where the branch energies $\Phi_{\mathrm{s}}$, $\Phi_{\mathrm{l}}$,
$\Phi_{\mathrm{r}}$, and $\Phi_{\mathrm{a}}$ are the repository's branchwise
quadratic corrections to \cref{eq:plasticity3d-elastic-branch}. The maintained
constitutive-AD path differentiates the smooth branches of
\cref{eq:plasticity3d-potential} directly with JAX rather than substituting a
hand-derived return-map tangent; the branch interfaces remain nonsmooth and are
treated as part of the algorithmic constitutive surrogate rather than as a
globally smooth continuum potential.

For the strain visualizations in this paper we use the deviatoric norm
\begin{equation}
  \|\varepsilon_{\mathrm{dev}}\|
  =
  \sqrt{
    \varepsilon^T
    \left(
      \operatorname{diag}(1,1,1,\tfrac{1}{2},\tfrac{1}{2},\tfrac{1}{2})
      - \frac{1}{3}mm^T
    \right)
    \varepsilon
  },
  \qquad
  m = (1,1,1,0,0,0)^T,
  \label{eq:plasticity3d-devnorm}
\end{equation}
which matches the engineering-shear convention of
\cref{eq:plasticity3d-engineering-strain}.

Because \cref{eq:plasticity3d-discrete,eq:plasticity3d-potential} define an
algorithmic endpoint surrogate rather than a path-consistent incremental
history functional, the results section reports this family with a separate
validation ladder. Same-case source-faithfulness to the source-style surrogate
implementation is treated separately from a fixed-lambda source-operator
comparison. The latter checks source-family behavior under matched geometry,
materials, boundary constraints, and reduction parameters, but it is not a
path-consistent incremental-history validation. The paper therefore keeps the
Plasticity3D claims at the level of a repository-specific surrogate
approximation rather than a path-consistent replacement. The related
convex-optimization and continuation literature
\citep{sysala2025convexoptimization,sysala2025advancedcontinuation} is used as
source-family context, not as evidence that this repository surrogate is a full
incremental elastoplastic solver.

\begin{figure}[H]
  \centering
  \includegraphics[width=\PaperFigureMedium]{plasticity3d_state_pair.pdf}
  \caption{Corrected glued-bottom Plasticity3D showcase at
  $\lambda_{\mathrm{sr}}=1.55$: displacement magnitude and surface deviatoric
  strain for the maintained-local \code{P4(L1\_2)} run, both derived from the
  scalar surrogate in
  \cref{eq:plasticity3d-discrete}.}
  \label{fig:plasticity3d-state}
\end{figure}

\begin{figure}[H]
  \centering
  \includegraphics[width=\PaperFigureMedium]{plasticity3d_convergence.pdf}
  \caption{Newton-history convergence for all nine completed corrected
  glued-bottom maintained-local Plasticity3D study runs at
  $\lambda_{\mathrm{sr}}=1.55$. The upper panel shows the value of the
  algorithmic surrogate from \cref{eq:plasticity3d-discrete}; the lower panel
  shows the gradient norm used as the nonlinear stopping metric, with color
  indicating degree and line style indicating mesh level.}
  \label{fig:plasticity3d-convergence}
\end{figure}

\begin{table}[H]
  \centering
  \begin{tabular*}{\textwidth}{@{}l@{\extracolsep{\fill}} c c c c@{}}
\toprule
Element & Free DOFs & Energy & $\|g\|_{\mathrm{final}}$ & Wall time [s] \\
\midrule
$P_{1}(L_{1})$ & \num{10526} & \num{-2953167} & $\num{4.56}\times 10^{-3}$ & \num{1.79} \\
$P_{1}(L_{2})$ & \num{79024} & \num{-2991497} & $\num{1.64}\times 10^{-3}$ & \num{8.72} \\
$P_{1}(L_{3})$ & \num{610964} & \num{-3004742} & $\num{4.07}\times 10^{-3}$ & \num{17.2} \\
$P_{1}(L_{4})$ & \num{4801816} & \num{-3009658} & $\num{1.16}\times 10^{-3}$ & \num{96.7} \\
$P_{2}(L_{1})$ & \num{79024} & \num{-3012813} & $\num{2.47}\times 10^{-3}$ & \num{18.3} \\
$P_{2}(L_{2})$ & \num{610964} & \num{-3013032} & $\num{2.96}\times 10^{-3}$ & \num{98.3} \\
$P_{2}(L_{3})$ & \num{4801816} & \num{-3013149} & $\num{1.14}\times 10^{-3}$ & \num{1330} \\
$P_{4}(L_{1})$ & \num{610964} & \num{-3013227} & $\num{5.53}\times 10^{-3}$ & \num{208} \\
$P_{4}(L_{2})$ & \num{4801816} & \num{-3013348} & $\num{4.42}\times 10^{-3}$ & \num{2298} \\
\bottomrule
\end{tabular*}

  \caption{All nine completed glued-bottom maintained-local Plasticity3D
  study rows at $\lambda_{\mathrm{sr}} = 1.55$, reported as endpoint values of
  \cref{eq:plasticity3d-discrete}.}
  \label{tab:plasticity3d-benchmark}
\end{table}

\begin{figure}[H]
  \centering
  \includegraphics[width=\PaperFigureFull]{plasticity3d_highest_mesh_y_slice_comparison.pdf}
  \caption{Shared-scale upper $y$-slice comparison of
  $\|\varepsilon_{\mathrm{dev}}\|$ from \cref{eq:plasticity3d-devnorm} at the
  highest mesh level of each degree line:
  \code{P1(L1\_2\_3\_4)}, \code{P2(L1\_2\_3)}, and \code{P4(L1\_2)}. The three
  runs use the corrected glued-bottom maintained-local solver stack at
  $\lambda_{\mathrm{sr}} = 1.55$ and are compared on one common color scale.}
  \label{fig:plasticity3d-highest-slice}
\end{figure}

Plasticity3D is the benchmark where detailed constitutive definitions are most
important: the paper later plots values of \cref{eq:plasticity3d-discrete} and
\cref{eq:plasticity3d-devnorm}, so both quantities are defined explicitly here.
Timing, scaling, and degree-versus-resolution evidence for this family are
reported in \Cref{sec:num-plasticity3d}.

\subsection{Topology Optimization}
\label{sec:bench-topology}

The literature model is the standard SIMP-style compliance topology-optimization
setting on a cantilever domain, together with filtering and continuation ideas
from the classical topology literature
\citep{bendsoe2003topology,sigmund2001topology,ferrari2020top99,bourdin2001filters}.
In the repository specialization studied here, the topology benchmark is a
reduced optimization problem rather than a single equilibrium minimization. The
design variable is a latent nodal field
$z_h \in Z_h$, mapped to a physical density through the logistic transform
\begin{equation}
  \theta_h(z_h)
  =
  \theta_{\min} + (1-\theta_{\min}) \sigma(z_h),
  \qquad
  \sigma(z) = \frac{1}{1+e^{-z}}.
  \label{eq:topology-density-map}
\end{equation}
For fixed density, the maintained discretization solves for the linear elastic
displacement state
$u_h(\theta_h)$ solving
\begin{equation}
  a_{\theta_h}(u_h,v_h)
  :=
  \sum_{e=1}^{n_{\mathrm{el}}}
    |e|\,\theta_e^{p}
    \varepsilon(u_h)^T C \varepsilon(v_h)
  =
  \ell(v_h)
  \qquad \forall v_h \in V_{h,0},
  \label{eq:topology-mechanics}
\end{equation}
on the cantilever domain $\Omega = (0,2)\times(0,1)$ with clamped left edge
and prescribed traction on a right-edge patch. The reduced objective solved in
the maintained code is
\begin{align}
  \mathcal{J}_h(z_h)
  &=
  \sum_{e=1}^{n_{\mathrm{el}}}
  |e|
  \Bigg[
    e_e^{\mathrm{frozen}} \theta_e^{-p}
    + \lambda_V \theta_e
    + \alpha_{\mathrm{reg}}
      \left(
        \frac{\ell_{\mathrm{pf}}}{2} \lvert \nabla \theta_e \rvert^2
        + \frac{1}{\ell_{\mathrm{pf}}} W(\theta_e)
      \right)
    + \frac{\mu_{\mathrm{move}}}{2} (z_e-z_e^{\mathrm{old}})^2
  \Bigg],
  \label{eq:topology-reduced-objective}
\end{align}
with $W(\theta) = \theta^2(1-\theta)^2$, SIMP-style continuation in $p$, and
repository regularization terms that play a mesh-control role analogous to
filtering/regularization in \citet{bourdin2001filters}. Here
$e_e^{\mathrm{frozen}}$ denotes the element elastic energy density from the
previous mechanics state; the latent logistic map, frozen-energy term, and move
penalty are repository-specific design choices rather than standard SIMP
notation. The maintained outer loop alternates mechanics solves for
\cref{eq:topology-mechanics} with reduced design updates for
\cref{eq:topology-reduced-objective}.

\begin{figure}[H]
  \centering
  \includegraphics[width=\PaperFigureMedium]{topology_density.pdf}
  \caption{Representative final density field for the maintained parallel
  topology benchmark.}
  \label{fig:topology-density}
\end{figure}

\begin{figure}[H]
  \centering
  \includegraphics[width=0.96\PaperFigureNarrow]{topology_history.pdf}
  \caption{Objective-history view for the reduced topology objective
  \cref{eq:topology-reduced-objective}, with compliance on the left axis and
  volume fraction plus the continuation schedule in $p$ on the right axis.}
  \label{fig:topology-history}
\end{figure}

\begin{table}[H]
  \centering
  \begin{tabular*}{\textwidth}{@{}l@{\extracolsep{\fill}} c c c c c@{}}
\toprule
Case & Ranks & Outer iters & Compliance & Volume fraction & Wall time [s] \\
\midrule
192$\times$96 & 1 & 121 & \num{4.1557} & \num{0.4} & \num{10} \\
768$\times$384 & 32 & 66 & \num{9.9074} & \num{0.3749} & \num{25.1} \\
\bottomrule
\end{tabular*}

  \caption{Representative topology endpoint values from
  \cref{eq:topology-reduced-objective}.}
  \label{tab:topology-benchmark}
\end{table}

For topology, the relevant notion of ``convergence'' is the reduced objective
history rather than a single equilibrium energy. Parallel timing and scaling
for the same family are reported in \Cref{sec:num-topology}.
